\chapter{总结与展望}

本文围绕 Ti/Sn 氧化物的可控生长机理及其耦合石墨烯后的电化学性能开展了较为系统的研究，对材料的制备过程、结构演化特征以及性能表现之间的关系进行了初步梳理。结合前文的实验结果与分析可以看出，Ti/Sn 氧化物体系在微观形貌调控、界面耦合优化以及电化学响应提升方面具有较好的研究潜力，所获得的结果为后续进一步完善相关材料体系提供了一定的实验依据。

从研究过程来看，阳极氧化条件、电解液环境以及后续复合过程都会对材料最终的形貌和性能产生明显影响。通过对不同实验参数下样品形貌与电化学特征的比较，可以认为，合理控制氧化物生长过程有助于改善材料的比表面积、电子传输路径以及活性位点分布，从而在一定程度上提升其储能性能。与此同时，石墨烯的引入为复合材料提供了更优的导电网络，使体系在电荷传输和结构稳定性方面表现出一定优势，这也说明复合化设计仍然是后续研究中值得继续关注的重要方向。

当然，本研究仍然存在一定的不足。例如，现有实验条件下对生长机理的讨论还主要建立在宏观现象与表征结果的基础之上，对于部分界面反应细节和结构演变过程的揭示还不够深入；同时，样品性能评价的维度仍有进一步扩展的空间，对循环稳定性、倍率性能及实际应用适配性的讨论还可以更加充分。由于实验周期和测试条件有限，部分参数尚未实现更精细的优化，这也使得本文结论更多体现为阶段性认识。

在后续工作中，可以从以下几个方面继续推进。首先，可进一步细化阳极氧化工艺参数的控制范围，结合更多原位或动态表征手段，对 Ti/Sn 氧化物的生长规律进行更深入的分析。其次，可围绕石墨烯与氧化物界面的结合方式开展优化设计，探索不同复合比例、不同组装方法对材料结构稳定性和导电性能的影响。最后，还可以结合实际器件构筑需求，对材料在更复杂工作条件下的电化学行为进行测试，以增强研究结果的应用价值。

总体而言，本文完成了对 Ti/Sn 氧化物相关体系的基础探索，并对其耦合石墨烯后的性能提升思路进行了初步验证。尽管研究内容仍有进一步深化的空间，但所得到的结果为后续相关研究提供了参考，也为高性能氧化物基电极材料的设计与优化积累了经验。

这里引用了一些文献来测试格式：关于刀具振动检测的研究见文献\cite{devillez2007tool}，热敏陶瓷的性能研究见文献\cite{chen2021microstructure}，智能刀具的综述见文献\cite{liu2021review}。此外，航空领域的经典著作有文献\cite{liu1988strake}，流体有限元法的学位论文见文献\cite{zhu1996new}，以及关于反馈声抵消器的专利文献\cite{li1986feedback}。
